\section*{ЗАКЛЮЧЕНИЕ}
\addcontentsline{toc}{section}{ЗАКЛЮЧЕНИЕ}

Прогресс в развитии нейронных сетей и методов машинного обучения предоставил новые возможности для создания систем биометрической идентификации. Современные технологии позволяют эффективно обрабатывать биометрические данные, такие как изображения лиц и отпечатков пальцев, обеспечивая высокую точность и надёжность аутентификации.

В условиях роста потребности в безопасных системах идентификации разработка интеллектуальных решений для распознавания пользователей является актуальной задачей. Применение таких систем позволяет повысить уровень безопасности, упростить процесс аутентификации и снизить затраты на реализацию.

Для решения задачи биометрической идентификации была разработана интеллектуальная система, основанная на использовании сверточной нейронной сети ResNet18 в составе сиамской архитектуры для распознавания отпечатков пальцев и модели InceptionResnetV1 для обработки лиц. Система реализована в виде настольного приложения с графическим интерфейсом.

Основные результаты работы::
\begin{enumerate}
	\item Реализован модуль захвата, обработки изображений лица и извлечения векторных признаков лиц.
	\item Реализован модуль распознавания отпечатков пальцев, реализованный на основе сиамской нейронной сети, позволяющей эффективно сравнивать изображения отпечатков.
	\item Графический интерфейс, разработанный с использованием библиотеки Tkinter, который обеспечивает интуитивно понятное взаимодействие с пользователем для добавления новых пользователей, создания QR-кодов и выполнения распознавания.
	\item Механизм обработки карт доступа, который интегрируется с процессом верификации для повышения безопасности.
\end{enumerate}

Все требования, объявленные в техническом задании, были полностью реализованы. Все задачи, поставленные в начале разработки проекта, были решены.

Готовый рабочий проект представлен в виде настольного приложения с графическим интерфейсом.